\section{总结与延伸}

\begin{figure}[H]
\centering
\includegraphics[width=0.95\textwidth]{slides-images/slide-065.jpg}
\caption{本节课总结\protect\footnotemark}
\end{figure}
\footnotetext{Slides 第66页。}

\subsection{核心要点回顾}

本节课完整覆盖了从预训练模型到对齐模型的后训练流程，核心内容包括：

\textbf{SFT 部分}：
\begin{enumerate}
    \item 指令微调数据有三种主要范式：NLP 任务聚合（FLAN）、LLM 生成（Alpaca）、人类众包（Open Assistant）
    \item ``高质量数据''的定义远比想象中复杂——完全正确的数据可能教模型幻觉
    \item 少量数据即可产生巨大影响（包括安全微调）
    \item 现代实践中，SFT 与预训练正通过中间训练阶段融合
\end{enumerate}

\textbf{RLHF 部分}：
\begin{enumerate}
    \item RLHF 数据收集面临诸多挑战：长度偏差、标注者偏见、事实核查困难、GPT-4 作弊
    \item AI 反馈（RLAIF）已成为主流趋势，效果可与人类标注媲美
    \item PPO 是经典但复杂的 RL 优化算法
    \item DPO 通过优雅的数学推导将 RL 问题化归为监督学习
    \item 警惕奖励过度优化和奖励黑客现象
\end{enumerate}

\subsection{开放问题与前沿方向}

\begin{itemize}
    \item \textbf{自我改进的上限在哪里？} 模型能从自身反馈中提取多少改进？理论上上限很高（模型包含了整个预训练语料的信息），但实际效果取决于具体实现
    \item \textbf{如何从根本上解决幻觉？} SFT 和 RLHF 都无法完全解决幻觉问题，可能需要在预训练阶段引入 RL 风格的自适应训练
    \item \textbf{可验证奖励的 RL}（下一讲主题）：在数学等可验证领域，可以用正确性作为奖励信号，避免人类/AI 标注的各种偏差
    \item \textbf{后训练的伦理问题}：标注工人的薪酬和工作条件、文化偏见的传播等
\end{itemize}

\subsection{拓展阅读}

\begin{itemize}
    \item Ouyang et al. (2022). \textit{Training language models to follow instructions with human feedback} (InstructGPT)
    \item Rafailov et al. (2023). \textit{Direct Preference Optimization: Your Language Model is Secretly a Reward Model} (DPO)
    \item Bai et al. (2022). \textit{Constitutional AI: Harmlessness from AI Feedback} (Anthropic)
    \item Schulman et al. (2017). \textit{Proximal Policy Optimization Algorithms} (PPO)
    \item Wang et al. (2023). \textit{How Far Can Camels Go? Exploring the State of Instruction Tuning on Open Resources}
    \item Hu et al. (2024). \textit{MiniCPM: Unveiling the Potential of Small Language Models with Scalable Training Strategies}
    \item Lambert et al. (2024). \textit{Tulu 3: Pushing Frontiers in Open Language Model Post-Training}
    \item Meng et al. (2024). \textit{SimPO: Simple Preference Optimization with a Reference-Free Reward}
    \item Zelikman et al. (2022). \textit{STaR: Bootstrapping Reasoning With Reasoning} 以及后续的 Quiet-STaR
    \item John Schulman. \textit{Reinforcement Learning from Human Feedback: Progress and Challenges} (Berkeley 报告)
\end{itemize}

\end{document}
